\documentclass[10pt,a4paper]{amsart}
\usepackage[margin=19mm]{geometry}
\usepackage{amsmath,amssymb,mathtools}
\usepackage[scr=rsfs,cal=euler]{mathalfa}
\usepackage{xfrac}
\usepackage[unicode,hidelinks,bookmarks=false]{hyperref}
\newcommand{\ie}{i.e.\ }
\newcommand{\cf}{cf.\ }
\newcommand{\resp}{respectively\ }
\newcommand{\Subsection}{\subsection}
\providecommand{\nobreakdash}{\nobreak\hbox{-}}
\setlength{\emergencystretch}{2em}
\multlinegap0pt
\usepackage{amssymb}
\usepackage{mathtools}
\usepackage{booktabs}
\newtheorem{proposition}{Proposition}
\title{Notes on the 33-point Erd\H{o}s--Szekeres problem}
\author{Bogdan Dumitru}
\date{Source: arXiv:2512.24061v1 (CC BY 4.0)}
\begin{document}
\maketitle

\noindent\emph{Source and licence.} Notes on the 33-point Erd\H{o}s--Szekeres problem. Version arXiv:2512.24061v1 (CC BY 4.0), distributed by arXiv under the Creative Commons Attribution 4.0 International licence, http://creativecommons.org/licenses/by/4.0/, as recorded in the arXiv licence metadata for this exact version and shown on the abstract page https://arxiv.org/abs/2512.24061v1. Full licence notice: \texttt{licenses/arxiv-2512.24061-CC-BY-4.0.txt}.\par\smallskip

\noindent\textit{Editorial scope.} Counted unit: the article's proposition prop:4criterion (source0.tex lines 114--117). The article's complete original proof of it is delivered with it (source0.tex lines 119--126, one \texttt{proof} environment of 8 lines). No mathematics is written, repaired or re-derived: the statement and the proof are the article's own lines, in the article's own notation, and no retained line is substituted or re-flowed.  No further source-local context is needed: every cross-reference of every retained line resolves inside the two delivered spans.  Nothing was omitted from any retained span, so the package carries no omission record.  No retained line contains a citation, so the wrapper carries no bibliography and no bibliography entry.  No retained line contains a figure environment, an image-inclusion command or drawing code, so no image is delivered and the package declares no asset dependency.  Editorial formatting only, in addition: an AMS \texttt{amsart} preamble that restates the article's own declarations and macros used by the retained lines, namely the environments \texttt{proposition} on separate counters, together with the standard packages those lines need.  The retained environments print in this document as Proposition 1 in order of appearance, whereas the article prints the same items with the numbering of its own full document.  The complete native source of the article as distributed in the arXiv e-print for this exact version is bundled unchanged as \texttt{sources/191849/source0.tex}.
\begin{proposition}[4-set criterion for convex position]\label{prop:4criterion}
Let $S$ be a finite set of points in the plane in general position.
Then $S$ is in convex position if and only if every 4-point subset of $S$ is in convex position.
\end{proposition}

\begin{proof}[Proof sketch]
If $S$ is in convex position, then every subset of $S$ is also in convex position, so every 4-point subset is convex.

Conversely, assume $S$ is \emph{not} in convex position. Then some point $p\in S$ is not a vertex of the convex hull of $S$, hence $p$ lies strictly inside the convex polygon $P=\mathrm{conv}(V)$, where $V\subseteq S$ is the set of hull vertices.
Triangulate $P$ by fixing a hull vertex $v_0\in V$ and drawing diagonals from $v_0$ to all non-adjacent hull vertices; this partitions $P$ into triangles whose vertices lie in $V$.
Since $p$ lies inside $P$, it lies inside one of these triangles, say $\triangle abc$ with $a,b,c\in V\subseteq S$.
Then the 4-set $\{a,b,c,p\}$ is not in convex position (one point lies inside the triangle formed by the other three), contradicting the assumption that every 4-point subset of $S$ is convex.
\end{proof}

\end{document}
